% Chapter 2 of The Replica Rosetta Stone: the random energy model.
\rstchapter{The Random Energy Model}{Replicas and the second moment}{replicas\dict the second moment}{ch:rem}
\stonefoot{proved}
\begin{refsection}[chapters/02-rem.bib]

\begin{chaptersummary}
Decode one of $2^N$ orthogonal codewords sent through Gaussian noise, and the posterior is a Gibbs measure on $2^N$ independent Gaussian energies: Derrida's random energy model, at an inverse temperature equal to the signal strength. Physics derives, by counting levels and by replicas, a free entropy that freezes at $\beta_c=\sqrt{2\log2}$, one-step replica-symmetry breaking with Parisi parameter $m=\beta_c/\beta$, and a decoding threshold at $\beta_c$, which is Shannon's limit $E_b/N_0=\ln2$. We prove all of it with first and second moments, a fractional moment at $n=m$ and Gaussian integration by parts, and compute $\lim N^{-1}\log\E Z_N^n$ at every real $n>0$: the one-step formula gives it exactly, and it is not analytic at $n=m$, which is why the integer moments continue to the wrong answer. Exact computations up to $N=1280$ and simulations up to $N=20$ show how slowly the limits are reached; what remains open concerns the replica method, not the model.
\end{chaptersummary}

\section{The question: one message among \texorpdfstring{$2^N$}{2\string^N}}\label{sec:rem:question}

A transmitter sends $N$ bits, a word $\sigma\in\{\pm1\}^N$, through Gaussian noise. For a random codebook, with which Shannon proved his coding theorem, the scores of the wrong codewords are independent, and the posterior is a measure on exponentially many configurations with independent random weights \parencite[ch.~6 and its notes]{MezardMontanari2009}. Here is a Gaussian version in which every step is exact.

\paragraph{Orthogonal codes.} Give message $\sigma$ the codeword $\lambda\sqrt N\,e_\sigma\in\R^{2^N}$, with $(e_\sigma)$ an orthonormal basis (as in pulse-position modulation), and add independent $\Normal(0,1)$ noise to every coordinate. If $\sigma^*$ is sent, the receiver observes the scores $y_\sigma=\lambda\sqrt N\,\ind\{\sigma=\sigma^*\}+X_\sigma$, with $X_\sigma$ i.i.d.\ $\Normal(0,1)$. Given $\sigma^*=\sigma$, the density of $y$ is the standard Gaussian density times $\exp(\lambda\sqrt N\,y_\sigma-\lambda^2N/2)$, so for a uniform prior the posterior is
\begin{equation}\label{eq:rem:posterior}
\Prob(\sigma^*=\sigma\mid y)=\frac{e^{\lambda\sqrt N\,y_\sigma}}{\sum_\tau e^{\lambda\sqrt N\,y_\tau}},
\end{equation}
a Gibbs measure on $2^N$ configurations whose scores, all but one, are independent standard Gaussians, at an inverse temperature equal to the signal strength, as on the Nishimori line of \cref{ch:nishimori}.\sidenote{Decoding with $\Prob(\sigma)\propto e^{\beta\sqrt N y_\sigma}$ at some other $\beta$ is the finite-temperature decoding of \textcite[ch.~6]{MezardMontanari2009}; \cref{ex:rem:mismatch} works out its phase diagram for orthogonal codes.} The energy is $\lambda^2$ per bit and the noise variance per dimension is $1=N_0/2$, so the signal-to-noise ratio per bit is $E_b/N_0=\lambda^2/2$; maximum a posteriori decoding picks the largest score. Without the planted score, what is left is the model that \textcite{Derrida1980,Derrida1981} introduced as the limit $p\to\infty$ of the $p$-spin glasses.

\begin{definition}[The random energy model]\label{def:rem:rem}
Let $(X_\sigma)_{\sigma\in\{\pm1\}^N}$ be i.i.d.\ $\Normal(0,1)$. The \term{random energy model}\index{random energy model} (REM) at inverse temperature $\beta\ge0$ has partition function and Gibbs measure
\[
Z_N(\beta)=\sum_{\sigma\in\{\pm1\}^N}e^{\beta\sqrt N X_\sigma},\qquad G_N(\sigma)=\frac{e^{\beta\sqrt N X_\sigma}}{Z_N(\beta)},
\]
and free entropy $\phi_N(\beta)=\frac1N\E\log Z_N(\beta)$. The \term{level} of $\sigma$ is $X_\sigma/\sqrt N$, minus its energy per spin. We write $\gibbs{f}=\sum_\sigma f(\sigma)G_N(\sigma)$ and $Y_N=\sum_\sigma G_N(\sigma)^2$ for the \term{participation ratio}\index{participation ratio}, the probability that two independent samples from $G_N$ (two replicas) coincide.
\end{definition}

No spins interact: the only structure is the number $2^N$ of levels and their variance $N$.\sidenote{In the language of \cref{ch:freeentropy}, the REM is the Gaussian disordered system with covariance $\E[\sqrt NX_\sigma\sqrt NX_\tau]=N\ind\{\sigma=\tau\}$, the extreme case in which distinct configurations are uncorrelated.} Throughout, $\beta_c=\sqrt{2\log2}$, so that $2^N=e^{N\beta_c^2/2}$, and $m=\beta_c/\beta$.

\begin{physicsbox}[Derrida 1980, 1981; Gross and M\'ezard 1984]
The free entropy of the REM is
\begin{equation}\label{eq:rem:phi}
\phi(\beta)=\begin{cases}\log2+\dfrac{\beta^2}2,&\beta\le\beta_c,\\[1.2ex]\beta\beta_c,&\beta\ge\beta_c,\end{cases}\qquad\beta_c=\sqrt{2\log2}.
\end{equation}
At $\beta_c$ the model freezes: its entropy $\log2-\beta^2/2$ reaches zero and stays there, the energy per spin sticks at the ground-state value $-\beta_c$, and the Gibbs measure condenses on a few configurations. Below the critical temperature replica symmetry is broken at one step, with Parisi parameter $m=\beta_c/\beta$: two replicas sit in the same configuration with probability $1-m$ and are otherwise unrelated, $P(q)=m\,\delta_0+(1-m)\,\delta_1$.
\end{physicsbox}

For decoding the box predicts a threshold. At $\beta=\lambda$ the planted score contributes $e^{\lambda^2N}$ to the posterior's partition function and the others $e^{N\phi(\lambda)}$, and $\lambda^2>\phi(\lambda)$ exactly when $\lambda>\beta_c$: decoding should succeed when the REM at the decoder's temperature is frozen, that is, when $E_b/N_0>\ln2$. That is Shannon's limit for unlimited bandwidth: the capacity $W\log_2(1+P/(N_0W))$ bits per second \parencite{Shannon1949} increases to $P/(N_0\ln2)$ as the bandwidth $W$ grows, and at the rate $P/E_b$ this is $E_b/N_0\ge\ln2$, about $-1.59$ decibels.

How do physicists derive the box, by counting levels as Derrida did and by replicas as Gross and M\'ezard did? And what is proved: the free entropy, the meaning of the Parisi parameter, the decoding threshold, and what the exact replica moments $\E Z_N^n$ say about the replica trick?

\section{The physics: counting levels and replicating}\label{sec:rem:physics}

\subsection{Derrida's level count}\label{sec:rem:count}

\begin{heuristic}[Derrida's microcanonical solution]
Count the configurations at level about $u$. Each level falls near $u$ with probability about $e^{-Nu^2/2}$, so their expected number is about $e^{Ns(u)}$, with the \term{entropy}\index{entropy!of a level} $s(u)=\log2-u^2/2$, positive for $\abs u<\beta_c$ and negative beyond. Assume that the typical number of levels near $u$ equals its mean when the mean is exponentially large, and is zero when the mean is exponentially small. Then, by Laplace's method,
\[
Z_N(\beta)\approx\int_{-\beta_c}^{\beta_c}e^{N[s(u)+\beta u]}\dd u,\qquad\phi(\beta)=\max_{\abs u\le\beta_c}\bigl[s(u)+\beta u\bigr].
\]
The exponent is stationary at $u=\beta$. If $\beta\le\beta_c$, then $\phi=\log2+\beta^2/2$ and the Gibbs measure lives at level $\beta$, on about $e^{Ns(\beta)}$ configurations. If $\beta>\beta_c$, the maximum sits at the edge $u=\beta_c$, where $s=0$, and $\phi=\beta\beta_c$: the Gibbs measure lives on the deepest levels, $e^{o(N)}$ of them, and the entropy is zero.
\end{heuristic}

Two steps are uncontrolled: ``typical equals mean'' for the random counts, and Laplace's method on a sum whose terms fluctuate. The second derivative $\phi''$ jumps from $1$ to $0$ at $\beta_c$, a continuous transition; because the Gibbs measure condenses there, it is often called a random first-order transition \parencite[ch.~5]{MezardMontanari2009}.\sidenote{The specific heat per spin, $\beta^2\phi''(\beta)$, jumps from $\beta_c^2=2\log2$ to $0$; in Derrida's units it jumps from $J^2/(2T_c^2)$ to $0$.}

\subsection{Replicas: the integer moments}\label{sec:rem:replicas}

The replica method computes $\E Z_N^n$ for integer $n$, continues the answer to real $n$, and uses $\E\log Z_N=\lim_{n\to0}\frac1n\log\E Z_N^n$. For the REM the integer moments are an exercise in counting.

\begin{computation}{The replicated partition function of the REM}\label{comp:rem:replicas}
\step \emph{Replicate and average.} For a positive integer $n$, $Z_N^n=\sum_{\sigma^1,\dots,\sigma^n}\exp(\beta\sqrt N\sum_{a}X_{\sigma^a})$, and the variance of $\sum_aX_{\sigma^a}$ is $\sum_{a,b}Q_{ab}$, with the \term{overlap matrix}\index{overlap!in the REM} $Q_{ab}=\ind\{\sigma^a=\sigma^b\}$; so $\E Z_N^n=\sum_{\sigma^1,\dots,\sigma^n}\exp(\frac{\beta^2N}2\sum_{a,b}Q_{ab})$.
\step \emph{Group.} The matrix $Q$ is the equivalence relation of a partition $\pi$ of the replicas into blocks of equal configurations; the blocks receive distinct configurations, in $(2^N)_{\abs\pi}=2^N(2^N-1)\cdots(2^N-\abs\pi+1)$ ways, and $\sum_{a,b}Q_{ab}=\sum_{B\in\pi}\abs B^2$. Hence, exactly at every $N$,
\begin{equation}\label{eq:rem:replicated}
\E Z_N^n=\sum_\pi(2^N)_{\abs\pi}\exp\Bigl(\frac{\beta^2N}2\sum_{B\in\pi}\abs B^2\Bigr).
\end{equation}
\step \emph{Saddle point.} There are finitely many partitions, so $\frac1N\log\E Z_N^n\to\max_\pi\bigl[\abs\pi\log2+\frac{\beta^2}2\sum_B\abs B^2\bigr]$. With $k$ blocks, $\sum_B\abs B^2$ is largest for one block of size $n-k+1$ and $k-1$ singletons, a convex function of $k$, so the maximum is at $k=n$ (all replicas apart) or $k=1$ (all together):
\[
\lim_{N\to\infty}\frac1N\log\E Z_N^n=\max\Bigl(n\bigl(\log2+\tfrac{\beta^2}2\bigr),\ \log2+\tfrac{n^2\beta^2}2\Bigr).
\]
The branches cross at $n=m^2$; above $\beta_c$ every integer $n\ge2$ lies on the second.
\end{computation}

\paragraph{Replica symmetry and its crisis.} The two partitions of step~3 are the two \term{replica-symmetric}\index{replica symmetry!in the REM} saddle points, $Q_{ab}=\delta_{ab}$ and $Q_{ab}=1$. The first, continued to $n\to0$, gives the annealed free entropy $\log2+\beta^2/2$ at every temperature, with entropy $\phi-\beta\phi'=\log2-\beta^2/2$, negative for $\beta>\beta_c$, which no system with finitely many configurations can have: replica symmetry must fail below the critical temperature. The second, continued, gives $(\log2+n^2\beta^2/2)/n\to\infty$. For $n<1$ the two exchange roles, and the replica method prescribes the smaller one; \textcite[ch.~8]{MezardMontanari2009} stress that nothing within the method justifies this choice.

\paragraph{One step of replica-symmetry breaking.} Parisi's ansatz divides the $n$ replicas into $n/m'$ groups of $m'$, with equal configurations within a group and distinct ones across groups. Then $\abs\pi=n/m'$ and $\sum_B\abs B^2=nm'$, so
\begin{equation}\label{eq:rem:g1rsb}
g_{\mathrm{1RSB}}(n,m')=n\Bigl(\frac{\log2}{m'}+\frac{\beta^2m'}2\Bigr).
\end{equation}
Continued in $n$, the integer constraint $1\le m'\le n$ is read as $n\le m'\le1$, over which $g$ is minimized \parencite[ch.~8]{MezardMontanari2009}. The minimum over $m'>0$ is at $m'=\beta_c/\beta=m$, with $\frac1n\,g_{\mathrm{1RSB}}(n,m)=\beta\beta_c$; it lies in $(0,1]$ exactly when $\beta\ge\beta_c$, and for $\beta<\beta_c$ the constrained minimum is the replica-symmetric $m'=1$. Both phases of \eqref{eq:rem:phi} come out, with zero entropy below $T_c$, and so does the overlap distribution: the fraction of pairs of replicas in a common group is $\frac{(n/m)m(m-1)}{n(n-1)}\to1-m$ as $n\to0$. Writing $Y_N$ as the limit $n\to0$ of $Z_N^{n-2}\sum_\sigma e^{2\beta\sqrt NX_\sigma}$, the same bookkeeping gives $\E Y_N\to1-m$, the condensation formula, which first appeared in the replica computation of \textcite{GrossMezard1984} \parencite[notes to ch.~5]{MezardMontanari2009}.\index{replica symmetry breaking!one step}\index{Parisi parameter}

\subsection{Replicas at finite \texorpdfstring{$n$}{n}}\label{sec:rem:finiten}

The one-step formula makes sense at every real $n>0$, not only as $n\to0$. With Parisi's rule, minimize over $m'\in[n,1]$ when $n<1$ and maximize over $m'\in[1,n]$ when $n>1$:
\begin{equation}\label{eq:rem:phi1rsb}
\varphi_{\mathrm{1RSB}}(n)=\begin{cases}\min_{n\le m'\le1}g_{\mathrm{1RSB}}(n,m'),&0<n\le1,\\ \max_{1\le m'\le n}g_{\mathrm{1RSB}}(n,m'),&n\ge1.\end{cases}
\end{equation}
Since $g_{\mathrm{1RSB}}(n,\cdot)$ is convex with its minimum at $m'=m$, for $\beta>\beta_c$ this is $n\beta\beta_c$ (at $m'=m$) when $n\le m$ and $\log2+n^2\beta^2/2$ (at $m'=n$, one group holding all the replicas) when $m<n\le1$; for $n\ge1$ the maximum sits at an endpoint and reproduces \cref{comp:rem:replicas}. The quadratic branch describes rare samples, in which one level sits at the atypical depth $u=n\beta>\beta_c$, an event of probability about $e^{-N(u^2/2-\log2)}$; \textcite[chs.~5, 8]{MezardMontanari2009} read $\lim N^{-1}\log\E Z_N^n$ as a Legendre transform of the large deviations of the free energy.

That is the failure of the naive replica trick: above $\beta_c$ every integer moment lies on the quadratic branch, whose interpolation $\log2+n^2\beta^2/2$ gives $\varphi(n)/n\to\infty$ as $n\to0$. Integer moments could not fix the continuation anyway: Carleman's criterion fails, since $\E Z_N^n\ge e^{N\beta^2n^2/2}$ \parencite[ch.~8]{MezardMontanari2009}, and adding $a\sin(\pi n)$ changes no integer value \parencite{Tanaka2007}.\sidenote{\Textcite{Tanaka2007} goes further and exhibits distributions with the same integer moments and different means of the logarithm, among them a two-spin random Ising model.} One-step breaking selects one continuation, \eqref{eq:rem:phi1rsb}, justified by its consequences: positive entropy, the right free entropy at both ends, a participation ratio in $[0,1]$. The uncontrolled steps are: typical equals mean; the continuation in $n$; the ansatz; the minimum over $[n,1]$; and the exchange of $n\to0$ with $N\to\infty$.

\section{The stone}\label{sec:rem:stone}

\begin{stone}{level counting and replicas}{the moment method}
\stonestep{Count levels}{About $e^{Ns(u)}$ levels near $u$ when $s(u)>0$, none when $s(u)<0$; Laplace's method over $u$.}{First moment: $\Prob(K_N(u)\ge1)\le\frac12e^{Ns(u)}$ removes the levels with $s<0$. Second moment: $\Var K_N(u)\le\E K_N(u)$, and Chebyshev makes the count typical when $s>0$ (\cref{lem:rem:levels}).}
\stonestep{Replicate}{$\E Z_N^n$ is a sum over partitions of the replicas, \eqref{eq:rem:replicated}; the saddle point picks one partition.}{Only $n=1$ and $n=2$ appear, as inequalities: Jensen's bound by $\E Z_N$, and the mean and variance of the level count.}
\stonestep{Replica symmetry and its crisis}{All replicas apart: the annealed free entropy $\log2+\beta^2/2$, whose entropy $\log2-\beta^2/2$ is negative below $T_c$, impossible with $2^N$ states.}{Jensen: $\phi_N\le\log2+\beta^2/2$, with equality in the limit when $\beta\le\beta_c$ (\cref{thm:rem:freeentropy}); no level lies above $\beta_c\sqrt N$, and the Gibbs entropy tends to $\max(\log2-\beta^2/2,0)$ (\cref{cor:rem:condensation}).}
\stonestep{One-step breaking}{Groups of $m'$ replicas; minimizing $n(\frac{\log2}{m'}+\frac{\beta^2m'}2)$ over $m'\in[n,1]$ gives $m=\beta_c/\beta$ and $\phi=\beta\beta_c$.}{A fractional moment: $\phi_N\le\frac1{Nm}\log\E Z_N^m$, and truncation at the extreme level $\beta_c\sqrt N$ gives $\E Z_N^m\le\frac32e^{Nm\beta\beta_c}$; the lower bound counts levels just below $\beta_c$.}
\stonestep{Parisi parameter}{$P(q)=m\delta_0+(1-m)\delta_1$; the weights fluctuate from sample to sample.}{$\E Y_N=1-\phi_N'(\beta)/\beta\to1-m$ by Gaussian integration by parts and convexity (\cref{cor:rem:condensation}); the law $PD(m,0)$ of the weights is a theorem (\cref{thm:rem:pd}).}
\stonestep{$n\to0$}{Continue in $n$ and exchange $n\to0$ with $N\to\infty$.}{No continuation: the exact $\varphi(n)$ (\cref{thm:rem:moments}) is linear up to $n=m$ and not analytic there, so the integers fix the wrong analytic function; the limits do commute.}
\end{stone}

The proof follows Derrida's count, not the replicas, and the first and second moments justify the count exactly: for independent levels, ``typical equals mean'' is Chebyshev's inequality and ``no level where the mean is small'' is Markov's. Where the physics extremizes over the group size $m'$, the proof optimizes a fractional moment,\index{fractional moment} and the optimal exponent is the Parisi parameter, the largest $n$ at which the bound $\phi_N\le\frac1{Nn}\log\E Z_N^n$ is sharp in the limit; where the physics reads $1-m$ off the replica matrix, the proof reads it off the derivative of the free entropy.\sidenote{In the SK model of \cref{ch:sk} the same exponent appears in Guerra's broken-symmetry bound, as the power $m$ in terms of the form $\frac1m\log\E(\cdot)^m$, and Ruelle's probability cascades take the place of the Poisson--Dirichlet weights.} One physical step has no counterpart: the proof never continues in $n$. It computes $\lim N^{-1}\log\E Z_N^n$ at every real $n$ (\cref{thm:rem:moments}), and the answer explains both the failure and the success of the replica trick: the integers determine the wrong analytic function, because the true one is singular at $n=m$, but the one-step formula with Parisi's rule is exact at every $n$.

\section{The proof: first and second moments}\label{sec:rem:proof}

\subsection{Gaussian tails and level counts}\label{sec:rem:tools}

\begin{lemma}[Gaussian tails]\label{lem:rem:gauss}
Let $X\sim\Normal(0,1)$, and write $H(x)=\Prob(X>x)$ throughout.
\begin{enumerate}
\item For $x\ge0$, $H(x)\le\frac12e^{-x^2/2}$; for $x>0$, $H(x)\ge\frac{x}{1+x^2}\,\frac{e^{-x^2/2}}{\sqrt{2\pi}}$.
\item For real $t$ and $a$, $\E[e^{tX}\ind\{X\le a\}]=e^{t^2/2}\bigl(1-H(a-t)\bigr)$ and $\E[e^{tX}\ind\{X>a\}]=e^{t^2/2}H(a-t)$.
\end{enumerate}
\end{lemma}

\begin{proof}
For the upper bound of part~1, substitute $y=x+v$ and use $e^{-xv}\le1$:
\[
H(x)=\int_0^\infty e^{-(x+v)^2/2}\frac{\dd v}{\sqrt{2\pi}}\le e^{-x^2/2}\int_0^\infty e^{-v^2/2}\frac{\dd v}{\sqrt{2\pi}}=\frac12e^{-x^2/2}.
\]
For the lower bound, $\frac{\dd}{\dd y}\bigl(-y^{-1}e^{-y^2/2}\bigr)=(1+y^{-2})e^{-y^2/2}$, so $x^{-1}e^{-x^2/2}=\int_x^\infty(1+y^{-2})e^{-y^2/2}\dd y\le(1+x^{-2})\sqrt{2\pi}\,H(x)$. Part~2 follows from $e^{tx-x^2/2}=e^{t^2/2}e^{-(x-t)^2/2}$, integrated over $x\le a$ and over $x>a$.
\end{proof}

For a level $u$ let $K_N(u)=\#\{\sigma:X_\sigma>u\sqrt N\}$, the number of configurations above level $u$, and $\mu_N(u)=\E K_N(u)=2^NH(u\sqrt N)$ its mean; recall Derrida's entropy $s(u)=\log2-u^2/2$.\index{level count}

\begin{lemma}[First and second moments of the level counts]\label{lem:rem:levels}
For every $N\ge1$ and real $u$:
\begin{enumerate}
\item (first moment) $\Prob(K_N(u)\ge1)\le\mu_N(u)$, and $\mu_N(u)\le\frac12e^{Ns(u)}$ if $u\ge0$;
\item (second moment) $\Var K_N(u)\le\mu_N(u)$, hence $\Prob\bigl(K_N(u)\le\frac12\mu_N(u)\bigr)\le4/\mu_N(u)$;
\item if $u>0$, then $\frac1N\log\mu_N(u)\ge s(u)-\frac1N\log\bigl(\sqrt{2\pi}\,(u\sqrt N+u^{-1}N^{-1/2})\bigr)$, so that $\frac1N\log\mu_N(u)\to s(u)$;
\item $N^{-1/2}\max_\sigma X_\sigma\to\beta_c$ almost surely.
\end{enumerate}
\end{lemma}

\begin{proof}
The plan is to treat $K_N(u)$ as a binomial count. It is a sum of $2^N$ independent indicators, each equal to $1$ with probability $q=H(u\sqrt N)$, so its mean is $2^Nq$ and its variance $2^Nq(1-q)\le\mu_N(u)$. Part~1 is Markov's inequality, $\Prob(K\ge1)\le\E K$, with the upper bound of \cref{lem:rem:gauss} and $2^N=e^{N\log2}$. Part~2 is Chebyshev's inequality: $\Prob(K\le\mu/2)\le\Prob(\abs{K-\mu}\ge\mu/2)\le4\Var K/\mu^2\le4/\mu$. Part~3 is the lower bound of \cref{lem:rem:gauss}; with part~1 it gives the limit. For part~4 fix $0<\eps<\beta_c$. By part~1, $\Prob(\max_\sigma X_\sigma>(\beta_c+\eps)\sqrt N)=\Prob(K_N(\beta_c+\eps)\ge1)\le\frac12e^{Ns(\beta_c+\eps)}$, and $s(\beta_c+\eps)<0$. By part~2, $\Prob(\max_\sigma X_\sigma\le(\beta_c-\eps)\sqrt N)=\Prob(K_N(\beta_c-\eps)=0)\le4/\mu_N(\beta_c-\eps)$, which decays exponentially by part~3 because $s(\beta_c-\eps)>0$. Both bounds are summable in $N$, and the Borel--Cantelli lemma gives the almost sure limit.
\end{proof}

\subsection{The free entropy}\label{sec:rem:fe}

\begin{theorem}[The free entropy of the REM]\label{thm:rem:freeentropy}
For every $\beta\ge0$,
\[
\lim_{N\to\infty}\frac1N\log Z_N(\beta)=\phi(\beta)=\max_{0\le u\le\beta_c}\bigl[s(u)+\beta u\bigr]\quad\text{almost surely and in }L^1,
\]
with $\phi$ as in \eqref{eq:rem:phi}; in particular $\phi_N(\beta)\to\phi(\beta)$. Moreover, for every $N\ge1$, $\phi_N(\beta)\le\log2+\beta^2/2$, and if $\beta>\beta_c$, then
\begin{equation}\label{eq:rem:upper}
\phi_N(\beta)\le\beta\beta_c+\frac{\log(3/2)}{mN},\qquad m=\frac{\beta_c}\beta.
\end{equation}
\end{theorem}

\begin{proof}
The plan is to reduce the claim to the mean $\phi_N$ by concentration, to bound the mean from above by Jensen's inequality, applied to $Z_N$ below $\beta_c$ and to the fractional power $Z_N^m$ above it, and to bound it from below by counting levels. If $\beta=0$, then $Z_N=2^N$ and there is nothing to prove; let $\beta>0$.

\emph{Step 1: concentration.} Let $F(x)=\log\sum_\sigma e^{\beta\sqrt Nx_\sigma}$ for $x\in\R^{2^N}$. For every $y$ and $\sigma$, $e^{\beta\sqrt Nx_\sigma}\le e^{\beta\sqrt N\norm{x-y}_\infty}e^{\beta\sqrt Ny_\sigma}$; summing over $\sigma$ and taking logarithms gives $\abs{F(x)-F(y)}\le\beta\sqrt N\norm{x-y}_\infty\le\beta\sqrt N\norm{x-y}_2$. Gaussian concentration for Lipschitz functions (\cref{thm:toolkit:lipschitz}) gives an absolute constant $c>0$ ($c=\frac12$ with the sharp constant) such that
\[
\Prob\Bigl(\Bigl|\frac1N\log Z_N-\phi_N\Bigr|\ge t\Bigr)\le2\exp\Bigl(-\frac{cNt^2}{\beta^2}\Bigr)\qquad(t>0).
\]
The right side is summable in $N$, so $\frac1N\log Z_N-\phi_N\to0$ almost surely by the Borel--Cantelli lemma, and integrating the tail gives $\E\abs{\frac1N\log Z_N-\phi_N}\le\beta\sqrt{\pi/(cN)}$. It remains to show that $\phi_N\to\phi$.

\emph{Step 2: upper bounds.} By Jensen's inequality, $N\phi_N\le\log\E Z_N=N(\log2+\beta^2/2)$, which is $N\phi(\beta)$ when $\beta\le\beta_c$. Let $\beta>\beta_c$. For $0<n\le1$, concavity of the logarithm gives $\E\log Z_N=\frac1n\E\log Z_N^n\le\frac1n\log\E Z_N^n$; we take $n=m$. Split $Z_N=Z_\le+Z_>$ according to whether $X_\sigma\le\beta_c\sqrt N$ or $X_\sigma>\beta_c\sqrt N$. The map $x\mapsto x^m$ is concave and subadditive on $[0,\infty)$, so
\[
\E Z_N^m\le\E Z_\le^m+\E Z_>^m\le(\E Z_\le)^m+\sum_\sigma\E\bigl[e^{m\beta\sqrt NX_\sigma}\ind\{X_\sigma>\beta_c\sqrt N\}\bigr].
\]
By part~2 of \cref{lem:rem:gauss} with $t=\beta\sqrt N$ and $a=\beta_c\sqrt N$, then part~1, and $2^N=e^{N\beta_c^2/2}$,
\[
\E Z_\le=2^Ne^{N\beta^2/2}H\bigl((\beta-\beta_c)\sqrt N\bigr)\le\tfrac12e^{N[\beta_c^2+\beta^2-(\beta-\beta_c)^2]/2}=\tfrac12e^{N\beta\beta_c}.
\]
With $t=m\beta\sqrt N=\beta_c\sqrt N$ each term of the sum equals $e^{N\beta_c^2/2}H(0)=\frac12e^{N\beta_c^2/2}$, and the sum is $\frac12e^{N\beta_c^2}=\frac12e^{Nm\beta\beta_c}$. Hence $\E Z_N^m\le(2^{-m}+\frac12)e^{Nm\beta\beta_c}\le\frac32e^{Nm\beta\beta_c}$, and $\phi_N\le\frac1{Nm}\log\E Z_N^m$ gives \eqref{eq:rem:upper}.

\emph{Step 3: lower bound.} Fix $u\in(0,\beta_c)$, write $\mu=\mu_N(u)$, and let $A=\{K_N(u)>\mu/2\}$. On $A$ more than $\mu/2$ configurations have $\beta\sqrt NX_\sigma>N\beta u$, so $\log Z_N\ge\log(\mu/2)+N\beta u$. On every sample, Jensen's inequality for the uniform average over $\sigma$ gives $\log Z_N\ge N\log2+\beta\sqrt N\,\bar X$, where $\bar X=2^{-N}\sum_\sigma X_\sigma$ is $\Normal(0,2^{-N})$. For $N$ so large that $\log(\mu/2)+N\beta u\ge0$, part~2 of \cref{lem:rem:levels} gives $\Prob(A)\ge1-4/\mu$, and therefore
\[
\E\log Z_N\ge\Bigl(1-\frac4\mu\Bigr)\Bigl(\log\frac\mu2+N\beta u\Bigr)-\beta\sqrt N\,\E\abs{\bar X}.
\]
By parts~1 and~3 of \cref{lem:rem:levels}, $\frac1N\log\mu\to s(u)>0$, so $\mu\to\infty$ exponentially fast; and $\E\abs{\bar X}\le2^{-N/2}$. Dividing by $N$ gives $\liminf_N\phi_N\ge s(u)+\beta u$ for every $u\in(0,\beta_c)$. The supremum of the right side over these $u$ is $\phi(\beta)$: it is attained at $u=\beta$ if $\beta<\beta_c$ and approached as $u\uparrow\beta_c$ if $\beta\ge\beta_c$.

\emph{Step 4: conclusion.} Steps 2 and 3 give $\phi_N\to\phi(\beta)$ for every $\beta$, and Step 1 transfers the limit to $\frac1N\log Z_N$, almost surely and in $L^1$.
\end{proof}

\begin{remark}[Other numbers of levels]\label{rem:rem:levels}
With $M_N$ levels, $2^{N-1}\le M_N\le2^N$, the proofs of \cref{thm:rem:freeentropy} and of part~4 of \cref{lem:rem:levels} go through verbatim, with $\mu_N(u)$ replaced by $M_NH(u\sqrt N)\ge\frac12\mu_N(u)$. We use this for $M_N=2^N-1$.
\end{remark}

\subsection{Energy, entropy and condensation}\label{sec:rem:condensation}

\begin{lemma}[Derivatives of convex limits]\label{lem:rem:convex}
Let $f_N$ be convex and differentiable on an open interval $I$, and $f_N\to f$ pointwise on $I$. If $f$ is differentiable at $\beta\in I$, then $f_N'(\beta)\to f'(\beta)$.
\end{lemma}

\begin{proof}
For $h>0$ with $\beta\pm h\in I$, convexity gives $\frac{f_N(\beta)-f_N(\beta-h)}h\le f_N'(\beta)\le\frac{f_N(\beta+h)-f_N(\beta)}h$. Let $N\to\infty$ and then $h\downarrow0$.
\end{proof}

\begin{corollary}[Energy, entropy and participation ratio]\label{cor:rem:condensation}
For every $\beta>0$, as $N\to\infty$,
\begin{enumerate}
\item $\frac1{\sqrt N}\E\gibbs{X_\sigma}\to\min(\beta,\beta_c)$: the Gibbs measure lives at level $\min(\beta,\beta_c)$;
\item $\frac1N\E S(G_N)\to\max\bigl(\log2-\frac{\beta^2}2,0\bigr)$, where $S(G_N)=-\sum_\sigma G_N(\sigma)\log G_N(\sigma)$;
\item $\E Y_N\to\max\bigl(1-\frac{\beta_c}\beta,0\bigr)$.
\end{enumerate}
\end{corollary}

\begin{proof}
The plan is to read all three statements off $\phi_N'(\beta)$, which converges by convexity, after computing that derivative a second time by Gaussian integration by parts. For each sample, $\beta\mapsto\log Z_N(\beta)$ is the logarithm of a sum of exponentials of linear functions of $\beta$, hence convex, and its derivative $\sqrt N\gibbs{X_\sigma}$ is bounded by the integrable $\sqrt N\max_\sigma\abs{X_\sigma}$. So $\phi_N$ is convex and, by dominated convergence, differentiable with $\phi_N'(\beta)=N^{-1/2}\E\gibbs{X_\sigma}$. The limit $\phi$ is differentiable on $(0,\infty)$ with $\phi'(\beta)=\min(\beta,\beta_c)$, since at $\beta_c$ both one-sided derivatives equal $\beta_c$. \Cref{lem:rem:convex,thm:rem:freeentropy} give part~1. For part~2, $-\log G_N(\sigma)=\log Z_N-\beta\sqrt NX_\sigma$, so $S(G_N)=\log Z_N-\beta\sqrt N\gibbs{X_\sigma}$ and $\frac1N\E S(G_N)=\phi_N-\beta\phi_N'\to\phi-\beta\phi'$, which is $\log2-\beta^2/2$ for $\beta\le\beta_c$ and $0$ above. For part~3, $G_N(\sigma)$ is a smooth function of the Gaussian vector $(X_\tau)$ with the bounded derivative $\partial G_N(\sigma)/\partial X_\sigma=\beta\sqrt N\bigl(G_N(\sigma)-G_N(\sigma)^2\bigr)$, so Gaussian integration by parts (\cref{thm:toolkit:gaussianibp}) gives
\[
\E\gibbs{X_\sigma}=\sum_\sigma\E\bigl[X_\sigma G_N(\sigma)\bigr]=\beta\sqrt N\sum_\sigma\E\bigl[G_N(\sigma)-G_N(\sigma)^2\bigr]=\beta\sqrt N\,(1-\E Y_N).
\]
Thus $\phi_N'(\beta)=\beta(1-\E Y_N)$, and $\E Y_N=1-\phi_N'(\beta)/\beta\to1-\min(\beta,\beta_c)/\beta$.
\end{proof}

Part~2 is the entropy crisis as a theorem, and part~3 says that the physicist's $P(1)=1-m$ is the mean participation ratio.\sidenote{In terms of the spin-glass overlap $R_{12}=\frac1N\ip{\sigma^1}{\sigma^2}$ of two samples, part~3 reads $\E\gibbs{R_{12}^2}\to\max(1-\beta_c/\beta,0)$: distinct configurations among the deepest levels are, by symmetry, nearly orthogonal (\cref{ex:rem:hamming}).} It does not say that $Y_N$ concentrates, and $Y_N$ does not: in the limit it is random, with the law that Ruelle's theorem describes (\cref{thm:rem:pd}).

\subsection{Decoding orthogonal codes}\label{sec:rem:decoding}

\begin{corollary}[The decoding threshold]\label{cor:rem:decoding}
Let $\sigma^*$ be uniform on $\{\pm1\}^N$, let $y_\sigma=\lambda\sqrt N\,\ind\{\sigma=\sigma^*\}+X_\sigma$ with $\lambda>0$, and let $\hat\sigma=\argmax_\sigma y_\sigma$ be the maximum a posteriori decoder.
\begin{enumerate}
\item If $\lambda>\beta_c$, then $\Prob(\hat\sigma=\sigma^*)\to1$ and $\Prob(\sigma^*\mid y)\to1$ in probability; if $\lambda<\beta_c$, both tend to $0$.
\item The mutual information satisfies $\frac1N I(\sigma^*;y)\to\min(\lambda^2/2,\log2)$.
\end{enumerate}
\end{corollary}

\begin{proof}
The plan is to split the posterior's partition function $Z(y)=\sum_\sigma e^{\lambda\sqrt Ny_\sigma}$ into the planted term $e^{\lambda\sqrt Ny_{\sigma^*}}$ and a REM with $2^N-1$ levels, $Z'=\sum_{\sigma\ne\sigma^*}e^{\lambda\sqrt NX_\sigma}$. By \cref{rem:rem:levels}, $N^{-1/2}\max_{\sigma\ne\sigma^*}X_\sigma\to\beta_c$ and $\frac1N\log Z'\to\phi(\lambda)$ almost surely, the second also in $L^1$; and $N^{-1/2}y_{\sigma^*}=\lambda+N^{-1/2}X_{\sigma^*}\to\lambda$.

\emph{Part~1.} Decoding succeeds exactly when $y_{\sigma^*}>\max_{\sigma\ne\sigma^*}X_\sigma$; divided by $\sqrt N$, the two sides tend to $\lambda$ and $\beta_c$. By \eqref{eq:rem:posterior}, $\Prob(\sigma^*\mid y)=\bigl(1+Z'e^{-\lambda\sqrt Ny_{\sigma^*}}\bigr)^{-1}$, and $\frac1N\log\bigl(Z'e^{-\lambda\sqrt Ny_{\sigma^*}}\bigr)\to\phi(\lambda)-\lambda^2$, which is $\lambda\beta_c-\lambda^2<0$ if $\lambda>\beta_c$ and $\log2-\lambda^2/2>0$ if $\lambda<\beta_c$.

\emph{Part~2.} The density of $y$ given $\sigma$ divided by its average over the uniform prior is $e^{\lambda\sqrt Ny_\sigma}/(2^{-N}Z(y))$, and $\E[\lambda\sqrt Ny_{\sigma^*}]=\lambda^2N$, so
\[
I(\sigma^*;y)=\E\log\frac{2^Ne^{\lambda\sqrt Ny_{\sigma^*}}}{Z(y)}=N(\log2+\lambda^2)-\E\log Z(y).
\]
With $a=\lambda\sqrt Ny_{\sigma^*}$ and $b=\log Z'$, $\max(a,b)\le\log Z(y)\le\log2+\max(a,b)$. Now $a/N\to\lambda^2$ and $b/N\to\phi(\lambda)$ in $L^1$, and $\abs{\max(a,b)-\max(a',b')}\le\abs{a-a'}+\abs{b-b'}$, so $\frac1N\E\log Z(y)\to\max(\lambda^2,\phi(\lambda))$. For $\lambda\ge\beta_c$ the maximum is $\lambda^2$ and $I/N\to\log2$; for $\lambda<\beta_c$ it is $\log2+\lambda^2/2$ and $I/N\to\lambda^2/2$.
\end{proof}

Below the threshold the observation carries $\lambda^2/2$ nats per bit, the whole energy budget, and above it all $N$ bits: orthogonal codes reach the capacity of the channel without bandwidth constraint.\index{orthogonal code}\index{Shannon limit}\sidenote{The cost is the bandwidth: $N$ bits use $2^N$ dimensions. Decoding itself is a scan of $2^N$ scores, linear in the size of the received signal; for random binary codes of rate $R$ the same scan takes time exponential in the block length.} The two decoding phases are the two phases of the REM at $\beta=\lambda$: above $\beta_c$ the planted level is deeper than every other and the frozen posterior sits on it; below, the $e^{N(\log2-\lambda^2/2)}$ wrong codewords at level $\lambda$ outweigh it.

\section{The proof: the replica moments, exactly}\label{sec:rem:moments}

The replica computation of \cref{sec:rem:finiten} makes a prediction at every real $n>0$. For the REM it can be checked.

\begin{theorem}[The replica moments of the REM]\label{thm:rem:moments}
Let $\beta>0$, $m=\beta_c/\beta$ and $n_c=\max(m,m^2)$. For every real $n>0$ the limit $\varphi(n)=\lim_{N\to\infty}\frac1N\log\E Z_N(\beta)^n$ exists, and
\begin{equation}\label{eq:rem:moments}
\varphi(n)=\begin{cases}n\,\phi(\beta),&0<n\le n_c,\\[0.6ex] \log2+\dfrac{n^2\beta^2}2,&n\ge n_c.\end{cases}
\end{equation}
\end{theorem}

For $n>1$ the proof compares $\E Z_N^n$ with $\E Z_N^{n-1}$.

\begin{lemma}[A recursion for the moments]\label{lem:rem:recursion}
For every real $n\ge1$ and every $N$,
\[
\E Z_N^n\le c_n\bigl(2^Ne^{Nn^2\beta^2/2}+\E Z_N\cdot\E Z_N^{n-1}\bigr),\qquad c_n=\max(1,2^{n-2}).
\]
\end{lemma}

\begin{proof}
Write $w_\sigma=e^{\beta\sqrt NX_\sigma}$ and $Z_{-\sigma}=Z_N-w_\sigma$, and let $\nu=n-1\ge0$. For $a,b\ge0$, $(a+b)^\nu\le c_n(a^\nu+b^\nu)$: if $\nu\le1$ by subadditivity, and if $\nu\ge1$ by convexity, $(\frac{a+b}2)^\nu\le\frac12(a^\nu+b^\nu)$. Hence
\[
Z_N^n=\sum_\sigma w_\sigma(w_\sigma+Z_{-\sigma})^\nu\le c_n\sum_\sigma\bigl(w_\sigma^n+w_\sigma Z_{-\sigma}^\nu\bigr).
\]
Since $w_\sigma$ is independent of $Z_{-\sigma}$ and $Z_{-\sigma}\le Z_N$, $\E[w_\sigma Z_{-\sigma}^\nu]\le\E w_\sigma\,\E Z_N^\nu$. Summing over $\sigma$, with $\E w_\sigma^n=e^{Nn^2\beta^2/2}$ and $\sum_\sigma\E w_\sigma=\E Z_N$, gives the claim.
\end{proof}

\begin{proof}[Proof of \cref{thm:rem:moments}]
The plan is to prove matching bounds on $\varphi_N(n)=\frac1N\log\E Z_N^n$ at exponential scale: lower bounds from typical samples and from samples with one deep level, upper bounds by the truncation of \cref{thm:rem:freeentropy} for $0<n\le1$ and by induction from \cref{lem:rem:recursion} for $n>1$. Write $\phi_a=\log2+\beta^2/2$ for the annealed free entropy, $f(n)=\log2+n^2\beta^2/2$, and $\varphi(n)$ for the right side of \eqref{eq:rem:moments}. Since $m^2\beta^2=\beta_c^2=2\log2$,
\begin{equation}\label{eq:rem:identity}
f(n)-n\phi_a=\frac{\beta^2}2(n-1)(n-m^2).
\end{equation}

\emph{Step 1: lower bounds.} Typical samples first: for $\eps>0$ the event $A_\eps=\{\frac1N\log Z_N\ge\phi-\eps\}$ has probability tending to $1$ by \cref{thm:rem:freeentropy}, and $\E Z_N^n\ge e^{nN(\phi-\eps)}\Prob(A_\eps)$, so $\liminf_N\varphi_N(n)\ge n\phi(\beta)$. Then one deep level: for $u\ge\beta_c$, $Z_N^n\ge e^{nN\beta u}\ind\{K_N(u)\ge1\}$. The levels are independent, so $\Prob(K_N(u)=0)=(1-H(u\sqrt N))^{2^N}\le e^{-\mu}$ with $\mu=\mu_N(u)$, and since $\mu\le\frac12$ by part~1 of \cref{lem:rem:levels}, $\Prob(K_N(u)\ge1)\ge1-e^{-\mu}\ge\mu-\frac12\mu^2\ge\frac12\mu$. With part~3 of \cref{lem:rem:levels}, $\liminf_N\varphi_N(n)\ge n\beta u+s(u)$, and the choice $u=\max(n\beta,\beta_c)$ gives $\Lambda(n)$, equal to $f(n)$ if $n\beta\ge\beta_c$ and to $n\beta\beta_c$ otherwise. So $\liminf_N\varphi_N(n)\ge\max(n\phi(\beta),\Lambda(n))$, and this maximum is $\varphi(n)$: if $\beta>\beta_c$, then $\Lambda(n)\ge n\beta\beta_c=n\phi$ and $\Lambda=\varphi$; if $\beta\le\beta_c$, then $\phi=\phi_a$, and \eqref{eq:rem:identity} shows that $\max(n\phi_a,\Lambda(n))$ is $n\phi_a$ up to $n=m^2$ and $f(n)$ beyond.

\emph{Step 2: upper bounds for $0<n\le1$.} If $\beta\le\beta_c$, Jensen's inequality for the concave map $x\mapsto x^n$ gives $\E Z_N^n\le(\E Z_N)^n=e^{Nn\phi_a}$. If $\beta>\beta_c$ and $n\beta\ge\beta_c$, subadditivity gives $\E Z_N^n\le\sum_\sigma\E e^{n\beta\sqrt NX_\sigma}=e^{Nf(n)}$. If $\beta>\beta_c$ and $n\beta<\beta_c$, the truncation of Step~2 of the proof of \cref{thm:rem:freeentropy}, with \cref{lem:rem:gauss}, gives
\begin{align*}
\E Z_N^n&\le(\E Z_\le)^n+2^Ne^{Nn^2\beta^2/2}H\bigl((\beta_c-n\beta)\sqrt N\bigr)\\
&\le2^{-n}e^{Nn\beta\beta_c}+\tfrac12e^{N[\beta_c^2+n^2\beta^2-(\beta_c-n\beta)^2]/2}\le\tfrac32e^{Nn\beta\beta_c}.
\end{align*}
In each case $\limsup_N\varphi_N(n)\le\varphi(n)$.

\emph{Step 3: $n>1$.} For $n>1$, \eqref{eq:rem:identity} shows that $\varphi(n)=\max(n\phi_a,f(n))$. Suppose that $\limsup_N\varphi_N\le\varphi$ on $(0,k-1]$ for an integer $k\ge2$, which Step~2 gives for $k=2$, and let $n\in(k-1,k]$ and $\nu=n-1$. By \cref{lem:rem:recursion} and $\log(a+b)\le\log2+\max(\log a,\log b)$,
\[
\limsup_N\varphi_N(n)\le\max\bigl(f(n),\ \phi_a+\varphi(\nu)\bigr).
\]
If $\varphi(\nu)=\nu\phi(\beta)$, then $\phi_a+\varphi(\nu)\le n\phi_a\le\varphi(n)$, because $\phi\le\phi_a$. Otherwise $\nu\ge n_c$ and $\varphi(\nu)=f(\nu)$, and $\phi_a+f(\nu)=f(n)+\log2-\nu\beta^2<f(n)$, because $\nu\beta^2\ge n_c\beta^2=\max(\beta\beta_c,\beta_c^2)\ge2\log2$. In both cases $\limsup_N\varphi_N(n)\le\varphi(n)$, which completes the induction; with Step~1 the theorem follows.
\end{proof}

\begin{corollary}[What the moments say about replicas]\label{cor:rem:replicas}
Let $\beta>0$.
\begin{enumerate}
\item $\varphi(n)/n=\phi(\beta)$ for $0<n\le n_c$. In particular $\lim_{n\downarrow0}\varphi(n)/n=\lim_N\phi_N(\beta)$: the limits $n\to0$ and $N\to\infty$ commute, and the fractional-moment bound $\phi_N\le\frac1{Nn}\log\E Z_N^n$ is sharp in the limit for every $n\le n_c$.
\item If $\beta>\beta_c$, then $\varphi$ is continuously differentiable at $n=m$, and $\varphi''$ jumps there from $0$ to $\beta^2$. If $\beta<\beta_c$, then $\varphi'$ jumps at $n=m^2$ from $\log2+\beta^2/2$ to $2\log2$.
\item If $\beta>\beta_c$, then $\varphi(k)=\log2+k^2\beta^2/2$ at every integer $k\ge1$, and the entire function $f(n)=\log2+n^2\beta^2/2$ agrees with $\varphi$ on $[m,\infty)$ but not on $(0,m)$, where $f(n)/n\to\infty$ as $n\downarrow0$.
\item The one-step formula with Parisi's rule is exact: $\varphi=\varphi_{\mathrm{1RSB}}$ of \eqref{eq:rem:phi1rsb} on $(0,\infty)$.
\end{enumerate}
\end{corollary}

\begin{proof}
Parts~1--3 are read off \eqref{eq:rem:moments}. At $n_c$ the left derivative is $\phi(\beta)$ and the right derivative $n_c\beta^2$: these are $\beta\beta_c$ and $m\beta^2=\beta\beta_c$ when $\beta>\beta_c$, and $\log2+\beta^2/2$ and $m^2\beta^2=2\log2$ when $\beta<\beta_c$. For part~4 let $h(m')=\frac{\log2}{m'}+\frac{\beta^2m'}2$, a convex function on $(0,\infty)$, decreasing up to $m'=m$ and increasing after, with $nh(1)=n\phi_a$, $nh(n)=f(n)$ and $nh(m)=n\beta\beta_c$. For $0<n\le1$ the minimum of $h$ over $[n,1]$ is at the point of $[n,1]$ nearest to $m$: at $m$ if $n\le m\le1$, giving $n\beta\beta_c$; at $n$ if $m<n$, giving $f(n)$; at $1$ if $m>1$, giving $n\phi_a$. These are the three cases of \eqref{eq:rem:moments} with $n\le1$. For $n\ge1$ the maximum of the convex $h$ over $[1,n]$ is at an endpoint, so $\varphi_{\mathrm{1RSB}}(n)=\max(n\phi_a,f(n))=\varphi(n)$.
\end{proof}

So the order of limits is harmless (part~1; \cref{ex:rem:commute} extends this to every Gaussian disordered system), and the continuation from the integers fails above $\beta_c$ (part~3) because the moments are singular at $n=m$, a phase transition in the replica number (part~2).\sidenote{\Textcite{OgureKabashima2009b} locate this singularity through the zeros of $\E Z_N^n$ in the complex $n$-plane: as $N$ grows they accumulate near the transition point, at an angle that marks a second-order transition, in agreement with part~2.} The one-step formula with Parisi's minimum is right at every $n$ (part~4), although nothing inside the replica method says why. Below $\beta_c$ the branch through $n=1$, all replicas apart, is the one that reaches $n\to0$, which is why replica symmetry is right there.\index{replica moments}

\subsection{What the general theory adds}\label{sec:rem:general}

Three further results complete the picture; we state them without proof.

\begin{theorem}[Poisson--Dirichlet weights; \textcite{Ruelle1987,BovierKurkova2004}]\label{thm:rem:pd}
Let $\beta>\beta_c$ and $m=\beta_c/\beta$. The point process $\sum_\sigma\delta_{G_N(\sigma)}$ of Gibbs weights converges in law to the point process of the weights $\xi_i/\sum_j\xi_j$, where $\xi_1>\xi_2>\cdots$ are the points of a Poisson process on $(0,\infty)$ with intensity $x^{-1-m}\dd x$. The law of these weights, in decreasing order, is the \term{Poisson--Dirichlet distribution}\index{Poisson--Dirichlet distribution} $PD(m,0)$.
\end{theorem}

\Textcite{Ruelle1987} reformulated Derrida's models with such Poisson processes, and \textcite{BovierKurkova2004} proved the convergence for the generalized REM, of which the REM is the case of one level. In size-biased order the weights are $W_i=V_i\prod_{j<i}(1-V_j)$ with independent $V_i\sim\mathrm{Beta}(1-m,im)$ \parencite{PitmanYor1997}. Since $\sum_{G_N(\sigma)<\eps}G_N(\sigma)^2\le\eps$, the theorem gives $Y_N\to\sum_iW_i^2$ in law and, $Y_N$ being bounded, in moments, so that $\Var Y_N\to\frac{m(1-m)}3>0$ (\cref{ex:rem:pdmoments}): the participation ratio does not self-average, as \textcite{DerridaToulouse1985} computed without replicas and \textcite{MezardParisiVirasoro1985} derived from independent, exponentially distributed free energies of the pure states.

\begin{theorem}[Fluctuations below $T_c$; \textcite{BovierKurkovaLowe2002}]\label{thm:rem:fluct}
Let $\beta>\beta_c$, $m=\beta_c/\beta$, $a_N=N\beta\beta_c-\frac1{2m}[\log(N\log2)+\log4\pi]$, and let $\mathcal P$ be a Poisson process on $\R$ with intensity $e^{-z}\dd z$. Then $\log Z_N-a_N\to\log\int e^{z/m}\mathcal P(\dd z)$ in law, and $\log Z_N-\E\log Z_N$ converges in law to the same limit, centered.
\end{theorem}

This is Theorem~1.5(ii) of the arXiv version of \textcite{BovierKurkovaLowe2002}, for our $Z_N$ rather than their $2^{-N}Z_N$; their Theorems~1.4 and~1.5(i) cover $\beta\le\beta_c$, with Gaussian fluctuations of $\log(Z_N/\E Z_N)$ for $\beta<\beta_c/2$ and non-Gaussian ones above. The two convergences differ by the deterministic shift $\E\log Z_N-a_N$, which therefore converges to the mean of the limit; with \cref{ex:rem:constant} this gives
\begin{equation}\label{eq:rem:finitesize}
\E\log Z_N=a_N+C_m+o(1),\qquad C_m=\frac{\log\Gamma(1-m)}m+\gamma\Bigl(\frac1m-1\Bigr),
\end{equation}
with $\gamma$ Euler's constant: $\phi_N$ approaches $\beta\beta_c$ from below like $\frac{\beta}{2\beta_c}\frac{\log N}N$.

Finally, Derrida's \term{generalized random energy model}\index{generalized random energy model} \parencite{Derrida1985} builds the energies as sums of independent contributions along the branches of a tree, and freezes in several steps. \Textcite{DerridaGardner1986} solved it, and \textcite{BovierKurkova2004}, who review the earlier rigorous work on its free energy, identified the limiting Gibbs measures with Ruelle's probability cascades, which reappear in Guerra's bound for the SK model (\cref{ch:sk}).

\section{Measuring it}\label{sec:rem:measuring}

Four scripts simulate the REM and compute exact finite-$N$ expectations. Since $Z_N$ is a sum of $2^N$ independent terms, $\E e^{-tZ_N}=L(t)^{2^N}$ with $L(t)=\E\exp(-te^{\beta\sqrt NX})$, and
\begin{gather*}
\E\log Z_N=\int_0^\infty\bigl(e^{-t}-L(t)^{2^N}\bigr)\frac{\dd t}t,\\
\E Z_N^n=\frac n{\Gamma(1-n)}\int_0^\infty\bigl(1-L(t)^{2^N}\bigr)\frac{\dd t}{t^{n+1}}\quad(0<n<1),
\end{gather*}
and $\E Y_N$ similarly, by the representation of \textcite[ch.~5]{MezardMontanari2009}.\sidenote{Exact finite-size expressions for the moments go back to \textcite{OgureKabashima2004,OgureKabashima2009a}. The script \texttt{02-rem-laplace.py} evaluates these integrals by quadrature in logarithmic variables, so that $N=1280$, with $2^N\approx10^{385}$ levels, costs no more than $N=20$; its checks against simulation at $N=8$--$12$ agree within the $95\%$ intervals, and refining the quadrature changes no reported digit.} All intervals are $95\%$ intervals: Student-$t$ for means, Clopper--Pearson for probabilities, percentile bootstrap where stated.

\begin{figure}[t]
\begin{fullwidth}
\centering
\begin{tikzpicture}
\begin{groupplot}[group style={group size=2 by 1,horizontal sep=17mm},rplot,width=80mm,height=42mm]
\nextgroupplot[xlabel={inverse temperature $\beta$},ylabel={free entropy},xmin=0,xmax=3.05,ymin=0.6,ymax=3.75,
  title={(a) $\phi_N(\beta)$ against $\beta$},legend pos=south east,legend style={font=\scriptsize}]
\addplot[physpred] table[x=beta,y=limit] {figures/data/02-rem-freeentropy-curve.dat};
\addplot[guide] table[x=beta,y=annealed] {figures/data/02-rem-freeentropy-curve.dat};
\addplot[proofpred] table[x=beta,y=bound20] {figures/data/02-rem-freeentropy-curve.dat};
\addplot[stone,line width=0.8pt,no marks] table[x=beta,y=exact20] {figures/data/02-rem-freeentropy-curve.dat};
\addplot[stone!50,line width=0.8pt,no marks] table[x=beta,y=exact320] {figures/data/02-rem-freeentropy-curve.dat};
\addplot[measured,ebars] table[x=beta,y=mean,y error plus=ci,y error minus=ci] {figures/data/02-rem-freeentropy-sim.dat};
\legend{$\phi(\beta)$,annealed,proved bound $N=20$,exact $N=20$,exact $N=320$,simulated $N=20$}
\nextgroupplot[xmode=log,xlabel={$N$},ylabel={$N\bigl(\phi(2)-\phi_N(2)\bigr)$},xmin=8,xmax=1600,ymin=0,ymax=7,
  title={(b) the approach to the limit at $\beta=2$},legend pos=south east,legend style={font=\scriptsize},
  xtick={10,100,1000},xticklabels={$10$,$100$,$1000$}]
\addplot[proofpred] table[x=N,y=asym] {figures/data/02-rem-freeentropy-size.dat};
\addplot[stone,line width=0.8pt,no marks] table[x=N,y=gap] {figures/data/02-rem-freeentropy-size.dat};
\addplot[measured,ebars] table[x=N,y=gap,y error plus=ci,y error minus=ci] {figures/data/02-rem-freeentropy-sizesim.dat};
\legend{asymptote \eqref{eq:rem:finitesize},exact,simulated}
\end{groupplot}
\end{tikzpicture}
\end{fullwidth}
\caption{The free entropy. (a) Derrida's formula \eqref{eq:rem:phi} (carnelian), the annealed value (grey dashed), the proved upper bound of \cref{thm:rem:freeentropy} at $N=20$ (lapis dashed), exact values at $N=20$ and $320$ (granite), and the mean of $\frac1N\log Z_N$ over $150$ simulated samples at $N=20$ (seed $2001$). (b) The deficit $N(\phi-\phi_N)$ at $\beta=2$: exact values, simulations ($1000$ samples for $N\le16$, $400$ at $N=18$, those of (a) at $N=20$; seeds $2002+N$), and the asymptote \eqref{eq:rem:finitesize}. Script: \texttt{figures/scripts/02-rem-freeentropy.py}.}\label{fig:rem:freeentropy}
\end{figure}

\paragraph{The free entropy.} In \cref{fig:rem:freeentropy}(a) the nine simulated means at $N=20$ all lie within their intervals of the exact finite-$N$ values, and above $\beta_c$ both lie well below the limit: at $\beta=2$ the simulation gives $2.214\pm0.015$, the exact value is $2.210$, and the limit is $2.355$. Panel (b) shows why: the deficit grows like $\frac1{2m}\log N=0.849\log N$, as \cref{thm:rem:fluct} predicts. The exact deficit exceeds the asymptote \eqref{eq:rem:finitesize} by $0.229$ at $N=20$, $0.028$ at $N=320$ and $0.011$ at $N=1280$, a remainder that decays like $(\log N)^2/N$ (its ratio to $(\log N)^2/N$ is $0.27$ at both of the last two sizes). Simulations at $N\le20$ alone could not separate a $\log N/N$ correction from a $1/N$ one.

\begin{figure}[t]
\begin{fullwidth}
\centering
\begin{tikzpicture}
\begin{groupplot}[group style={group size=2 by 1,horizontal sep=17mm},rplot,width=80mm,height=42mm]
\nextgroupplot[xlabel={inverse temperature $\beta$},ylabel={$\E Y_N$},xmin=0.55,xmax=3.05,ymin=0,ymax=0.75,
  title={(a) the participation ratio},legend pos=south east,legend style={font=\scriptsize}]
\addplot[physpred] table[x=beta,y=limit] {figures/data/02-rem-gibbs-curve.dat};
\addplot[stone,line width=0.8pt,no marks] table[x=beta,y=exact20] {figures/data/02-rem-gibbs-curve.dat};
\addplot[stone!70,line width=0.8pt,densely dashed,no marks] table[x=beta,y=exact80] {figures/data/02-rem-gibbs-curve.dat};
\addplot[stone!45,line width=0.8pt,no marks] table[x=beta,y=exact320] {figures/data/02-rem-gibbs-curve.dat};
\addplot[measured,ebars] table[x=beta,y=mean,y error plus=ci,y error minus=ci] {figures/data/02-rem-gibbs-sim.dat};
\legend{{$\max(1-\beta_c/\beta,0)$},exact $N=20$,exact $N=80$,exact $N=320$,simulated $N=20$}
\nextgroupplot[xlabel={replica number $n$},ylabel={$\frac1{Nn}\log\E Z_N^n$},xmin=0,xmax=2.15,ymin=2.1,ymax=4.6,
  title={(b) the replica moments at $\beta=2$},legend style={at={(axis cs:0.42,4.55)},anchor=north west,font=\scriptsize}]
\addplot[physpred] table[x=n,y=limit] {figures/data/02-rem-moments-curve.dat};
\addplot[guide] table[x=n,y=integer] {figures/data/02-rem-moments-curve.dat};
\addplot[stone,line width=0.8pt,densely dotted,no marks,restrict x to domain=0:0.97] table[x=n,y=N12] {figures/data/02-rem-moments-exact.dat};
\addplot[stone,line width=0.8pt,no marks,restrict x to domain=0:0.97] table[x=n,y=N320] {figures/data/02-rem-moments-exact.dat};
\addplot[measured,ebars] table[x=n,y=est,y error plus=hi,y error minus=lo] {figures/data/02-rem-moments-sim.dat};
\addplot[measured3,restrict x to domain=0.99:2.5] table[x=n,y=N12] {figures/data/02-rem-moments-exact.dat};
\legend{$\varphi(n)/n$,integers continued,exact $N=12$,exact $N=320$,sampled $N=12$,exact integer moments}
\end{groupplot}
\end{tikzpicture}
\end{fullwidth}
\caption{Condensation and the replica moments. (a) The mean participation ratio: the limit of \cref{cor:rem:condensation} (carnelian), exact values at $N=20$, $80$, $320$ (granite), and the mean of $Y_{20}$ over $150$ simulated samples (seed $2101$). (b) At $\beta=2$ ($m=0.589$): the limit $\varphi(n)/n$ of \cref{thm:rem:moments}, which is also the one-step formula with Parisi's rule (carnelian); the continuation $\frac{\log2}n+\frac{n\beta^2}2$ of the integer moments (grey dashed); exact values at $N=12$ and $320$, and at $n=1,2$ (squares); and the empirical mean of $Z_{12}^n$ over $20\,000$ samples (seed $2201$), with percentile-bootstrap intervals ($400$ resamples). Scripts: \texttt{figures/scripts/02-rem-gibbs.py}, \texttt{02-rem-moments.py}.}\label{fig:rem:gibbs}
\end{figure}

\paragraph{Condensation.} In \cref{fig:rem:gibbs}(a) the six simulated means of $Y_{20}$ lie within their intervals of the exact values, which approach the limit slowly: at $\beta=2$ they are $0.466$, $0.429$ and $0.417$ at $N=20$, $80$ and $320$, against $1-m=0.411$, a deficit that shrinks roughly like $\log N/N$; at $\beta_c$, where the limit is $0$, they are $0.163$, $0.079$ and $0.039$. For the Poisson--Dirichlet law, $10^4$ stick-breaking samples of $PD(m,0)$ ($3000$ sticks, seed $2102$) give $\E\sum_iW_i^2$ and $\E(\sum_iW_i^2)^2$ within their intervals of $1-m$ and $(1-m)(3-2m)/3$ at $m=0.25$, $0.5$, $0.75$ (for instance $0.503\pm0.006$ and $0.337\pm0.006$ at $m=0.5$). At $N=20$ and $\beta=2$ the simulated $\E Y_{20}^2=0.274\pm0.049$ is consistent with the $PD$ value $0.300$ at the effective parameter $m_{\mathrm{eff}}=1-\E Y_{20}=0.534$, but $150$ samples cannot separate it from the limiting $0.250$ at $m=0.589$.

\paragraph{The replica moments.} \Cref{fig:rem:gibbs}(b) shows the singularity of \cref{thm:rem:moments} forming. At every finite $N$ the curve $n\mapsto\frac1{Nn}\log\E Z_N^n$ is smooth; at $N=320$ it lies within $0.016$ of $\beta\beta_c=2.355$ below $n=0.5$ and on the limit above $n=0.65$. The integer moments lie on the limit, and the curve that continues them leaves it at $n=m$. Sampling cannot measure the quadratic branch: the empirical moments at $N=12$ agree with the exact values up to $n=0.7$ and fall short beyond, $2.533$ against $2.693$ at $n=1$ and $2.857$ against $4.347$ at $n=2$, because for $n>m$ the expectation is carried by samples with a deep level, of probability about $10^{-8}$ at $n=1$, far below $1/20\,000$; the bootstrap intervals cannot see what the sample does not contain.

\paragraph{Decoding.} The script \texttt{02-rem-decoding.py} computes the success probability of maximum a posteriori decoding, $\E\Phi(\lambda\sqrt N+X)^{2^N-1}$, by quadrature and the mutual information by the planted Laplace representation. The threshold of \cref{cor:rem:decoding} emerges slowly: at $\lambda=1$ the success probability falls from $0.440$ at $N=12$ to $0.215$, $0.021$ and $0.000$ at $N=40$, $160$, $640$, and at $\lambda=1.5$ it rises from $0.932$ to $0.990$, $1.000$, $1.000$ (\cref{ex:rem:threshold} locates the transition). The information per bit at $\lambda=1$ is $0.459$, $0.489$ and $0.499$ at $N=12$, $32$, $96$, against the limit $\frac12$. At $N=12$, $2000$ simulated transmissions at each of eleven values of $\lambda$ (seed $2301$) put ten of the eleven information estimates and eight of the eleven success frequencies within their intervals of the exact values; at the three outlying $\lambda$, $40\,000$ further transmissions (seed $2302$) agree with them. No measurement contradicts a proved result.

\section{The status board}\label{sec:rem:status}

For the REM the prediction of Derrida, Gross and M\'ezard is a theorem in every part, down to the replica moments at every real $n$ with their one-step formula, and the fluctuations are known at every temperature. What is false is what the physics itself rejects: the replica-symmetric solution below $T_c$, and the naive continuation of the integer moments.

\begin{statusboard}
\statusrow{Free entropy $\phi(\beta)$, freezing at $\beta_c=\sqrt{2\log2}$, zero entropy below $T_c$}{\proved}{\cref{thm:rem:freeentropy,cor:rem:condensation}; \textcite[ch.~9]{Bovier2006}}
\statusrow{Replica-symmetric (annealed) solution below $T_c$}{\knownfalse}{its entropy $\log2-\beta^2/2$ is negative; \cref{cor:rem:condensation}}
\statusrow{One-step breaking: $\E Y_N\to1-\beta_c/\beta$, $P(q)=m\delta_0+(1-m)\delta_1$}{\proved}{\cref{cor:rem:condensation}; predicted by \textcite{GrossMezard1984}}
\statusrow{Poisson--Dirichlet weights $PD(m,0)$; $Y_N$ not self-averaging}{\proved}{\textcite{Ruelle1987,BovierKurkova2004}; predicted by \textcite{MezardParisiVirasoro1985,DerridaToulouse1985}}
\statusrow{Fluctuations of $\log Z_N$ at every $\beta$}{\proved}{\textcite{BovierKurkovaLowe2002}, Thms 1.4--1.5 (arXiv version)}
\statusrow{Replica moments $\varphi(n)$ at real $n>0$, given by the one-step formula}{\proved}{\cref{thm:rem:moments,cor:rem:replicas}; transition in $n$ found by \textcite{OgureKabashima2004,OgureKabashima2009a}}
\statusrow{Continuation of the integer moments to $n\to0$}{\knownfalse}{for every $\beta>\beta_c$; \cref{cor:rem:replicas}, part~3}
\statusrow{Orthogonal codes: threshold $\lambda=\beta_c$ ($E_b/N_0=\ln2$), information $\min(\lambda^2/2,\log2)$}{\proved}{\cref{cor:rem:decoding}}
\statusrow{Generalized REM: free entropy, Ruelle cascades}{\proved}{\textcite{BovierKurkova2004}; predicted by \textcite{DerridaGardner1986}}
\end{statusboard}

\section{Where the translation fails}\label{sec:rem:fails}

\paragraph{What was assumed.} Independence of the energies, first of all: in the $p$-spin glasses at finite $p$ the energies of nearby configurations are correlated, for $p=2$ (the SK model of \cref{ch:sk}) the physics predicts full replica-symmetry breaking, and the counting arguments do not survive. Second, Gaussian levels: the freezing and the Poisson--Dirichlet weights come from the Gaussian tail. In the decoding reading the independence is that of orthogonal codewords, which occupy $2^N$ dimensions; random binary codes also give independent scores \parencite[ch.~6]{MezardMontanari2009}, structured codes do not. Third, the decoder uses the Bayes posterior; at another temperature a glassy phase appears (\cref{ex:rem:mismatch}).

\paragraph{Where the physics is wrong, or not yet right.} The replica-symmetric solution is wrong below $T_c$, and the physics knows it; the naive continuation of the integer moments is wrong above $\beta_c$ (\cref{cor:rem:replicas}). The one-step formula with Parisi's rule is right at every $n$, but the rule itself, a minimum over $m'\in[n,1]$, has no derivation inside replica theory. Finite size is where the physical description is least accurate: at $N=20$ and $\beta=2$ the free entropy is $0.145$ below its limit and $\E Y_N$ is $0.055$ above it, and replicas reproduce such corrections only if the numbers and sizes of Parisi's blocks are given negative variances \parencite{DerridaMottishaw2015}.

\paragraph{Open problems.} For the REM itself the leading order and the fluctuations are settled; the open questions concern the method. First, for integer-valued energies one-step breaking gives the free energy but not the overlaps, and \textcite{DerridaMottishaw2022} suggest that the block sizes should fluctuate and take complex values, as for the finite-size corrections. A derivation of such a rule, or a class of models in which it is proved to work, is not known to us (as of 2026). Second, for which models is $\varphi(n)=\lim N^{-1}\log\E Z_N^n$ linear on an interval $(0,n_c]$ and given at every $n$ by Parisi's formula with the constraint $n\le m'\le1$? \Cref{thm:rem:moments} settles the REM, and \cref{ex:rem:grem} poses the generalized REM. Third, \textcite[notes to ch.~8]{MezardMontanari2009} observe that the replica method with broken symmetry still awaits a mathematical framework; the REM shows that it must explain why the extremum is a minimum over $[n,1]$.

\begin{exercises}
\item \consolidation (The plain second moment.) (a) Show from \eqref{eq:rem:replicated} that $\E Z_N^2/(\E Z_N)^2=1+2^{-N}(e^{N\beta^2}-1)$, which tends to $1$ if and only if $\beta<\beta_c/\sqrt2$, and deduce $\phi_N\to\log2+\beta^2/2$ there with the Paley--Zygmund inequality and Step~1 of the proof of \cref{thm:rem:freeentropy}. (b) Show that $\E\tilde Z_N^2/(\E\tilde Z_N)^2\to1$ for every $\beta<\beta_c$, where $\tilde Z_N$ sums over $X_\sigma\le\beta_c\sqrt N$: which samples spoil the plain second moment? \hint{for $\beta_c/2<\beta<\beta_c$ the excess is at most a constant times $e^{-N(\beta_c-\beta)^2}$.}\label{ex:rem:secondmoment}

\item \consolidation (Moments of $PD(m,0)$.) (a) With the size-biased pick $W_1=V_1\sim\mathrm{Beta}(1-m,m)$, show that $\E\sum_iW_i^r=\E W_1^{r-1}=\frac{\Gamma(r-m)}{\Gamma(r)\Gamma(1-m)}$. (b) Show that $\E(\sum_iW_i^2)^2=\frac{(1-m)(3-2m)}3$, so that $Y_N$ does not converge in probability when $\beta>\beta_c$ (\cref{thm:rem:pd}). \hint{$\E\sum_{i\ne j}W_i^2W_j^2=\frac{m(1-m)^2}6$ is the probability that, of four independent draws from the weights, the first two fall on one atom and the last two on another.}\label{ex:rem:pdmoments}

\item \consolidation (Negative moments.) Show that $\varphi(n)=n\phi(\beta)$ for every real $n<0$. \hint{Jensen's inequality gives the lower bound; for the upper bound combine $Z_N\ge2^Ne^{\beta\sqrt N\bar X}$ with the Chernoff bound $\Prob(K_N(u)\le\frac12\mu_N(u))\le e^{-\mu_N(u)/8}$.}\label{ex:rem:negative}

\item \byhand (The binary REM of random codes.) Give $2^{NR}$ messages independent codewords uniform on $\{0,1\}^N$ and send one over the binary symmetric channel with flip probability $p<\frac12$. (a) Show that the expected number of wrong codewords at Hamming distance about $N\delta$ from the output is $e^{N[(R-1)\log2+h(\delta)]+o(N)}$, with $h$ the binary entropy in nats, so that the closest lies at distance $N\delta_{\mathrm{GV}}+o(N)$, where $h(\delta_{\mathrm{GV}})=(1-R)\log2$ and $\delta_{\mathrm{GV}}<\frac12$, and that maximum-likelihood decoding succeeds if and only if $R<1-h(p)/\log2$, the capacity. (b) With posterior weights $(p/(1-p))^d$, show that the free entropy of the wrong codewords is frozen at $-\delta_{\mathrm{GV}}\log\frac{1-p}p$ if $p<\delta_{\mathrm{GV}}$ and equals $(R-1)\log2-\log(1-p)$ otherwise \parencite[ch.~6]{MezardMontanari2009}: as for orthogonal codes, the posterior is frozen exactly when decoding succeeds.\label{ex:rem:binary}

\item \byhand (Decoding at the wrong temperature.) Decode orthogonal codes with $G_\beta(\sigma)\propto e^{\beta\sqrt Ny_\sigma}$. (a) Show that $G_\beta(\sigma^*)\to1$ in probability if $\beta\lambda>\phi(\beta)$ and $\to0$ if $\beta\lambda<\phi(\beta)$: the threshold is $\lambda=\beta_c$ for $\beta\ge\beta_c$ and $\lambda=\frac{\beta^2+\beta_c^2}{2\beta}$ for $\beta<\beta_c$ ($1.266$ at $\beta=0.8$). (b) Draw the phase diagram in the $(\lambda,1/\beta)$ plane, with ordered, glassy and paramagnetic phases, and compare with that of \textcite[ch.~6]{MezardMontanari2009} for random binary codes. Where does the Bayes line $\beta=\lambda$ meet the boundary?\label{ex:rem:mismatch}

\item \byhand (The constant of the extreme-value asymptote.) Let $S=\int e^{z/m}\mathcal P(\dd z)$, with $\mathcal P$ a Poisson process of intensity $e^{-z}\dd z$ and $0<m<1$. Show that $\E e^{-tS}=\exp(-\Gamma(1-m)t^m)$ and, from $\log s=\int_0^\infty(e^{-t}-e^{-ts})\frac{\dd t}t$, that $\E\log S=\frac{\log\Gamma(1-m)}m+\gamma(\frac1m-1)$; at $\beta=2$ this is $1.708$. Check \eqref{eq:rem:finitesize} against \cref{fig:rem:freeentropy}(b): $6.217$ exact against $6.207$ at $N=1280$.\label{ex:rem:constant}

\item \translate (The Hamming overlap.) The physicist's overlap in the REM is $Q_{12}=\ind\{\sigma^1=\sigma^2\}$; the spin-glass overlap is $R_{12}=\frac1N\ip{\sigma^1}{\sigma^2}$. Using the exchangeability of $(X_\sigma)$, show that $\E\gibbs{R_{12}^2}=\E Y_N+(1-\E Y_N)\frac{2^N/N-1}{2^N-1}$, deduce $\E\gibbs{R_{12}^2}\to\max(1-\beta_c/\beta,0)$, and state $P(q)=m\delta_0+(1-m)\delta_1$ as a theorem about two samples from $G_N$.\label{ex:rem:hamming}

\item \translate (Why the limits commute.) Suppose that $\log Z_N$ is an $L\sqrt N$-Lipschitz function of a standard Gaussian vector and that every $L_F$-Lipschitz $F$ satisfies $\E e^{t(F-\E F)}\le e^{t^2L_F^2/2}$. Show that, for $n>0$, $\frac1N\log\E Z_N^n$ lies between $n\phi_N$ and $n\phi_N+n^2L^2/2$; deduce $\lim_{n\downarrow0}\varphi(n)/n=\lim_N\phi_N$ when both limits exist, and check against \eqref{eq:rem:moments} with $L=\beta$: what fails in the replica trick is the continuation, not the order of limits.\label{ex:rem:commute}

\item \measurement (The finite-size decoding threshold.) Compute the $\lambda$ of $\frac12$ success for $N=20$, $80$, $320$ (the script's function \texttt{map\_success} gives $1.089$, $1.148$, $1.168$). From the Gumbel limit of the largest wrong score, derive $\lambda_{1/2}=\beta_c-\frac{\log(N\log2)+\log4\pi-2\gamma}{2\beta_cN}+o(N^{-1})$, and compare. Show that $\sqrt N$ times the width from $10\%$ to $90\%$ success tends to $2\times1.2816$: the noise on the transmitted score, not the fluctuation of the wrong scores, sets the width.\label{ex:rem:threshold}

\item \research\harder (The replica moments of the generalized REM.) For Derrida's generalized REM with two levels \parencite{Derrida1985}, in the notation of \textcite[\S1.2]{Kistler2015}, compute $\varphi(n)=\lim N^{-1}\log\E Z_N^n$ for real $n>0$ by the methods of \cref{thm:rem:moments}, locate its singularities, and compare with the two-step replica formula with Parisi's rule and with the nested bound $\E\log Z_N\le\frac1{m_1}\log\E[(\E[Z_N^{m_2}\mid\text{first generation}])^{m_1/m_2}]$, the analogue of Guerra's bound of \cref{ch:sk}. \hint{write $Z_N=\sum_{\sigma_1}e^{\beta X^{(1)}_{\sigma_1}}Z^{(2)}_{\sigma_1}$, a sum over the first generation of independent REMs of the second.}\label{ex:rem:grem}
\end{exercises}

\notesandrefs
\Textcite{Derrida1980,Derrida1981} introduced the REM as the limit $p\to\infty$ of the $p$-spin glasses and solved it by the microcanonical count of \cref{sec:rem:count}; \textcite[notes to ch.~8]{MezardMontanari2009} credit the replica solution to the same papers. \Textcite{GrossMezard1984} treated the $p$-spin model at large $p$ with replicas and one-step breaking, \textcite{MezardParisiVirasoro1985} and \textcite{DerridaToulouse1985} found the sample-to-sample fluctuations of the weights, and \textcite{Derrida1985} and \textcite{DerridaGardner1986} introduced and solved the generalized REM.

On the rigorous side, \textcite{OlivieriPicco1984} studied the thermodynamic limit of the REM, and \textcite{Ruelle1987}, \textcite{BovierKurkovaLowe2002} and \textcite{BovierKurkova2004} proved the results of \cref{sec:rem:general}. Textbook accounts are \textcite[chs.~9--10]{Bovier2006} and the lectures of \textcite{Kistler2015}; \textcite[chs.~5, 6, 8]{MezardMontanari2009} give the physical treatment, prove the level count by first and second moments, and treat the random code ensemble as a REM. The orthogonal-code reading of \cref{sec:rem:question}, with the capacity formula of \textcite{Shannon1949}, is its Gaussian analogue.

\Textcite{OgureKabashima2004} found, from exact finite-size moments of a discrete REM, a phase transition at a real replica number, and \textcite{OgureKabashima2009a,OgureKabashima2009b} identified the loss of analyticity in $n$ as the origin of one-step breaking in the REM; \textcite{DerridaMottishaw2015} obtain one-step breaking without an ansatz from contour integrals over complex replica numbers, and \textcite{Tanaka2007} discusses the moment problem behind the replica method. The proofs of \cref{thm:rem:freeentropy,thm:rem:moments} and \cref{cor:rem:condensation,cor:rem:decoding,cor:rem:replicas} were written for this book (the fractional-moment bound at $n=m$, arranged to mirror the one-step computation, the recursion of \cref{lem:rem:recursion}, and the derivative argument for $\E Y_N$); their ingredients are standard.

\paragraph{Translating the literature.} \Textcite{MezardMontanari2009} use energies of variance $N/2$ and the weight $e^{-\beta E}$: their $\beta$ is $\sqrt2$ times ours, their $\beta_c=2\sqrt{\log2}$, their energy density $\eps$ is our $-u/\sqrt2$, their $s_a(\eps)=\log2-\eps^2$ our $s(u)$, their $m_s(\beta)$ our $m$, and their $\psi(n)$ our $\varphi(n)$. \Textcite{Derrida1981} uses energies of variance $NJ^2/2$, so that our $\beta$ is $J/(\sqrt2\,T)$ and $T_c=J/(2\sqrt{\log2})$. \Textcite{BovierKurkovaLowe2002} write $Z_{\beta,N}$ for our $2^{-N}Z_N$ and $\alpha=\beta/\sqrt{2\ln2}$ for our $1/m$. \Textcite{Kistler2015} gives the levels variance $N$: his $X_\sigma$ is our $\sqrt NX_\sigma$, and his $\mathcal N_N(\lambda)$, the number of levels at least $\lambda N$, our $K_N(\lambda)$. \Textcite{PitmanYor1997} write $PD(\alpha,\theta)$; ours is the case $\alpha=m$, $\theta=0$.

\printchapterbib
\end{refsection}
